% main.tex -- AI-Guided Prediction of Thermoelastic Wave
% Propagation in Geological Media. Diploma defence deck, IITU 2026.

\documentclass[aspectratio=169,10pt]{beamer}
\usetheme{default}

\usefonttheme{professionalfonts}
\usepackage[T1]{fontenc}
\usepackage{lmodern}
\usepackage{amsmath}
\usepackage{amssymb}
\usepackage{booktabs}
\usepackage{array}
\usepackage{graphicx}

\definecolor{AcademicNavy}{HTML}{17324D}
\definecolor{AcademicBlue}{HTML}{2E5E88}
\definecolor{AcademicText}{HTML}{202124}
\definecolor{AcademicMuted}{HTML}{5F6B76}
\definecolor{AcademicRule}{HTML}{CBD5E1}
\definecolor{AcademicPanel}{HTML}{F6F8FA}

% Clean academic defence style: light background, dark text, restrained accents.
\setbeamertemplate{navigation symbols}{}
\setbeamertemplate{headline}{}
\setbeamertemplate{blocks}[default]
\setbeamersize{text margin left=8mm,text margin right=8mm}

\setbeamercolor{background canvas}{bg=white}
\setbeamercolor{normal text}{fg=AcademicText,bg=white}
\setbeamercolor{frametitle}{fg=AcademicNavy,bg=white}
\setbeamercolor{title}{fg=AcademicNavy,bg=white}
\setbeamercolor{subtitle}{fg=AcademicMuted,bg=white}
\setbeamercolor{author}{fg=AcademicText,bg=white}
\setbeamercolor{institute}{fg=AcademicMuted,bg=white}
\setbeamercolor{date}{fg=AcademicMuted,bg=white}
\setbeamercolor{structure}{fg=AcademicBlue}
\setbeamercolor{block title}{fg=AcademicNavy,bg=AcademicPanel}
\setbeamercolor{block body}{fg=AcademicText,bg=AcademicPanel}
\setbeamercolor{block title example}{fg=AcademicNavy,bg=AcademicPanel}
\setbeamercolor{block body example}{fg=AcademicText,bg=AcademicPanel}
\setbeamercolor{block title alerted}{fg=AcademicNavy,bg=AcademicPanel}
\setbeamercolor{block body alerted}{fg=AcademicText,bg=AcademicPanel}
\setbeamercolor{footline}{fg=AcademicMuted,bg=white}

\setbeamerfont{title}{series=\bfseries,size=\Large}
\setbeamerfont{subtitle}{size=\normalsize}
\setbeamerfont{frametitle}{series=\bfseries,size=\large}
\setbeamerfont{normal text}{size=\normalsize}
\setbeamerfont{block title}{series=\bfseries,size=\small}
\setbeamerfont{block body}{size=\small}
\setbeamerfont{footline}{size=\scriptsize}

\setbeamertemplate{itemize item}{\raise0.2ex\hbox{\tiny$\blacktriangleright$}}
\setbeamertemplate{itemize subitem}{\raise0.2ex\hbox{\tiny$\circ$}}
\setbeamertemplate{enumerate item}{\textbf{\insertenumlabel.}}

\setbeamertemplate{frametitle}{%
  \vspace{2mm}
  \insertframetitle\par
  \vspace{1mm}
  {\color{AcademicRule}\rule{\linewidth}{0.45pt}}
  \vspace{-1mm}
}

\setbeamertemplate{footline}{%
  \leavevmode
  \hbox{%
    \begin{beamercolorbox}[wd=.72\paperwidth,ht=2.5ex,dp=1ex,leftskip=8mm]{footline}%
      \color{AcademicMuted}\insertshorttitle
    \end{beamercolorbox}%
    \begin{beamercolorbox}[wd=.28\paperwidth,ht=2.5ex,dp=1ex,right,rightskip=8mm]{footline}%
      \color{AcademicMuted}\insertframenumber\,/\,\inserttotalframenumber
    \end{beamercolorbox}%
  }%
}

\setbeamertemplate{title page}{%
  \vspace{10mm}
  {\usebeamerfont{title}\usebeamercolor[fg]{title}\inserttitle\par}
  \vspace{3mm}
  {\color{AcademicRule}\rule{0.68\linewidth}{0.6pt}\par}
  \vspace{4mm}
  {\usebeamerfont{subtitle}\usebeamercolor[fg]{subtitle}\insertsubtitle\par}
  \vfill
  {\usebeamerfont{author}\usebeamercolor[fg]{author}\insertauthor\par}
  \vspace{2mm}
  {\usebeamerfont{institute}\usebeamercolor[fg]{institute}\insertinstitute\par}
  \vspace{2mm}
  {\usebeamerfont{date}\usebeamercolor[fg]{date}\insertdate\par}
}

% -- Metadata -----------------------------------
\title[AI-Guided Thermoelastic Prediction]{%
  AI-Guided Prediction of Thermoelastic Wave Propagation%
  \texorpdfstring{\\}{ }%
  in Geological Media%
}
\subtitle{A comparative Beamer presentation for diploma defence}

\author[Group of 4]{%
  Kasymkhanov Timur \and Askarov Insar \and Kaidrakhmetova Dinara \and Nurshanov Dias%
}
\institute{International Information Technology University}
\date{Bachelor's diploma defence, 2026}

\begin{document}

\begin{frame}[plain]
  \titlepage
  \vfill
  \begin{center}\footnotesize
    Research adviser: \textbf{Bakhyt N. Alipova} (PhD, Applied Mathematics)
  \end{center}
\end{frame}

\section{Motivation and Problem}

\begin{frame}{Relevance of the Topic}
  \begin{block}{Why this topic matters}
    \begin{itemize}
      \item Thermally affected rocks are relevant to geothermal systems,
            underground engineering, and geophysical interpretation.
      \item Temperature change may alter stress, deformation, and wave behavior
            inside the geological medium.
      \item Repeated coupled simulations are informative, but expensive for
            comparative directional studies.
    \end{itemize}
  \end{block}
  \vskip 6pt
  \begin{exampleblock}{Working position of this thesis}
    The project is a \emph{research prototype} for
    \emph{AI-guided directional prediction}, not a calibrated
    field-scale geological simulator.
  \end{exampleblock}
\end{frame}

\begin{frame}{Research Problem}
  \begin{exampleblock}{Problem statement}
    How can we estimate the directional characteristics of thermoelastic
    wave propagation in geological media without rerunning the full coupled
    numerical simulation for every scenario?
  \end{exampleblock}
  \vskip 6pt
  \begin{itemize}
    \item Real rocks are heterogeneous, porous, fractured, and often anisotropic.
    \item A complete field-scale thermoelastic description is beyond the scope
          of this bachelor thesis.
    \item The task is therefore framed as \emph{comparative directional prediction}
          under simplified 2D effective-medium assumptions.
  \end{itemize}
\end{frame}

\begin{frame}{Aim and Tasks}
  \begin{block}{Aim}
    To develop and justify an AI-guided framework for directional
    prediction of thermoelastic wave propagation in geological media.
  \end{block}
  \vskip 4pt
  \begin{enumerate}
    \item Establish the physical and geological basis of the problem.
    \item Formulate a simplified thermoelastic prediction task in 2D.
    \item Build a COMSOL-derived synthetic dataset for controlled comparison.
    \item Compare PINN, FNO, MeshGraphNet, and Transformer-style surrogates.
    \item Interpret the results, limitations, and next development steps.
  \end{enumerate}
\end{frame}

\begin{frame}{Thermoelastic Coupling Logic}
  \begin{enumerate}
    \item \textbf{Thermal excitation:} a local source changes the temperature field.
    \item \textbf{Thermal expansion:} temperature change produces local strain.
    \item \textbf{Stress formation:} constrained expansion creates thermal stress.
    \item \textbf{Mechanical response:} stress gradients drive displacement.
    \item \textbf{Propagation:} the response travels through the medium.
    \item \textbf{picture:} picture here
  \end{enumerate}
\end{frame}

\section{Physical Basis}

\begin{frame}{Table - Geological Media Parameter}
  \tiny
  \renewcommand{\arraystretch}{1.08}
  \setlength{\tabcolsep}{1.6pt}
  \resizebox{\textwidth}{!}{%
  \begin{tabular}{@{}llrrrrrrrrrrrrrl@{}}
    \toprule
    \textbf{Material} & \textbf{Group} & $\rho$ & $E$ & $\nu$ & $\mu$ & $K$ &
    $\lambda$ & $\alpha$ & $k$ & $C_p$ & $C$ & $\gamma$ & $V_P$ & $V_S$ &
    \textbf{Porosity} \\
    & & kg/m$^3$ & GPa & -- & GPa & GPa & GPa & $10^{-6}$/K &
    W/mK & J/kgK & MJ/m$^3$K & MPa/K & m/s & m/s & \% \\
    \midrule
    Granite & igneous intrusive & 2650 & 76.28 & 0.245 & 30.63 & 49.84 & 29.42 & 7.9 & 2.50 & 850 & 2.25 & 1.18 & 5850 & 3400 & total 45.5 \\
    Granodiorite & igneous intrusive & 2795 & 85.91 & 0.255 & 34.24 & 58.35 & 35.52 & -- & 2.40 & -- & -- & -- & 6100 & 3500 & -- \\
    Basalt & igneous volcanic & 3000 & 103.96 & 0.266 & 41.07 & 73.95 & 46.57 & -- & 1.95 & 1300 & 3.90 & -- & 6550 & 3700 & total 19.0 \\
    Diabase & igneous hypabyssal & 2975 & 98.22 & 0.274 & 38.56 & 72.36 & 46.66 & -- & 2.50 & 850 & 2.53 & -- & 6450 & 3600 & -- \\
    Gabbro & igneous intrusive mafic & 3075 & 114.27 & 0.287 & 44.40 & 89.33 & 59.72 & -- & 2.30 & 1000 & 3.08 & -- & 6950 & 3800 & total 43.5 \\
    Sandstone & sedimentary siliciclastic & 2725 & 55.75 & 0.259 & 22.13 & 38.61 & 23.86 & 11.6 & 2.25 & 1000 & 2.73 & 1.34 & 5000 & 2850 & total 31.5; eff. 26.5 \\
    Limestone & sedimentary carbonate & 2675 & 61.48 & 0.277 & 24.08 & 45.90 & 29.85 & 8.0 & 2.55 & 1050 & 2.81 & 1.10 & 5400 & 3000 & total 31.5; eff. 18.0 \\
    Marble & metamorphic carbonate & 2725 & 82.41 & 0.270 & 32.43 & 59.82 & 38.20 & 9.8 & 2.80 & 850 & 2.32 & 1.76 & 6150 & 3450 & -- \\
    Schist & metamorphic foliated & 2800 & 68.20 & 0.267 & 26.91 & 48.82 & 30.88 & -- & 2.75 & 1150 & 3.22 & -- & 5500 & 3100 & total 26.5; eff. 27.5 \\
    Quartzite & metamorphic quartz-rich & 2650 & 83.50 & 0.250 & 33.40 & 55.70 & 33.44 & -- & 5.10 & 1000 & 2.65 & -- & 6150 & 3550 & -- \\
    \bottomrule
  \end{tabular}%
  }
  \vskip 5pt
  \scriptsize
  Values are effective midpoint parameters used for model-ready comparison.
  Elastic constants are derived from $\rho$, $V_P$, and $V_S$ under isotropic
  assumptions; missing thermal expansion values are shown as dashes.

  MAKE A ROWS
\end{frame}

\begin{frame}{Rock Properties Controlling Propagation}
  \begin{columns}[T]
    \begin{column}{0.48\textwidth}
      \begin{block}{Mechanical controls}
        \begin{itemize}
          \item Density controls inertia.
          \item Elastic stiffness controls resistance to deformation.
          \item Porosity, cracks, and saturation change effective stiffness.
        \end{itemize}
      \end{block}
    \end{column}
    \begin{column}{0.48\textwidth}
      \begin{block}{Thermal controls}
        \begin{itemize}
          \item Conductivity controls heat transfer.
          \item Heat capacity controls thermal storage.
          \item Thermal expansion converts temperature change into strain.
        \end{itemize}
      \end{block}
    \end{column}
  \end{columns}
  \vskip 6pt
  Directional behavior can appear because bedding, foliation, aligned
  pores, fractures, and stress state make real rocks structurally ordered.
\end{frame}

\begin{frame}{Elastic Waves in Geological Media}
  \begin{columns}[T]
    \begin{column}{0.48\textwidth}
      \begin{block}{P-waves}
        Compressional waves. Particle motion is mainly parallel to the
        propagation direction and involves volume change.
      \end{block}
    \end{column}
    \begin{column}{0.48\textwidth}
      \begin{block}{S-waves}
        Shear waves. Particle motion is mainly transverse to the direction
        of propagation and depends on shear rigidity.
      \end{block}
    \end{column}
  \end{columns}
  \vskip 6pt
  \[
    V_P = \sqrt{\frac{K+\frac{4}{3}G}{\rho}},
    \qquad
    V_S = \sqrt{\frac{G}{\rho}}
  \]
  A stiffer rock usually transmits elastic disturbances faster, while
  porosity and compliant cracks often reduce effective wave velocity.
\end{frame}

\begin{frame}{Heat Conduction and Thermal Expansion}
  \begin{block}{Thermal part of the process}
    \[
      q = -k\nabla T,
      \qquad
      \rho C_p\frac{\partial T}{\partial t}
      =
      \nabla\cdot(k\nabla T)+Q
    \]
  \end{block}
  \vskip 4pt
  \begin{block}{Mechanical effect of heating}
    \[
      \theta = T-T_0,
      \qquad
      \varepsilon_{\mathrm{th}}=\alpha\Delta T
    \]
  \end{block}
  \vskip 4pt
  Heat diffuses through the solid framework, while thermal expansion
  creates strain. If expansion is constrained, thermal stress appears.
\end{frame}

\begin{frame}{Key Physical Parameters}
  \begin{block}{Change the name og this formula}
    \[
      \alpha_T = \frac{k}{\rho C_p}
    \]
  \end{block}
  \vskip 6pt
  \begin{itemize}
    \item $\rho$ controls inertia in the equation of motion.
    \item $K$ and $G$ control compressional and shear stiffness.
    \item $k$ and $C_p$ control how quickly temperature disturbances spread.
    \item $\alpha$ controls how strongly temperature change produces strain.
  \end{itemize}
\end{frame}

\begin{frame}{Thermoelastic Coupling}
  \begin{block}{Hooke law}
    \[
      \theta = T - T_0
    \]
    \[
      \sigma_{ij} = \lambda\,\delta_{ij}\,\varepsilon_{kk}
      + 2\mu\,\varepsilon_{ij}
      - \gamma\,\delta_{ij}\,\theta,
      \qquad
      \gamma = 3K\alpha
    \]
    \[
      \rho\,\frac{\partial^2u_i}{\partial t^2}
      = \frac{\partial \sigma_{ij}}{\partial x_j} + f_i
    \]
  \end{block}
  \vskip 6pt
  These relations provide the theoretical backbone of the thesis and
  motivate the physics-aware interpretation of the surrogate models.
\end{frame}

\begin{frame}{Mathematical Scope and Assumptions}
  \begin{itemize}
    \item The geological medium is represented as a 2D continuum:
          \[
            \Omega\subset\mathbb{R}^{2},\qquad (x,t)\in\Omega\times[0,T_f].
          \]
    \item The displacement field is in-plane:
          \[
            u(x,t)=
            \begin{bmatrix}
              u(x,t)\\
              v(x,t)
            \end{bmatrix}.
          \]
    \item Small deformation and linear isotropic thermoelasticity are used
          as first-order approximations.
    \item Pores, fractures, fluids, plasticity, damage, and full anisotropic
          constitutive modelling are not resolved explicitly.
  \end{itemize}
\end{frame}

\begin{frame}{Primary Fields in the 2D Model}
  \footnotesize
  \begin{tabular}{@{}lll@{}}
    \toprule
    \textbf{Field} & \textbf{Meaning} & \textbf{Role in the thesis} \\
    \midrule
    $T(x,t)$ & Temperature & Thermal state of the rock \\
    $\theta(x,t)=T-T_0$ & Temperature perturbation & Drives thermal strain/stress \\
    $(u,v)$ & In-plane displacement & Mechanical response at the probe \\
    $\varepsilon_{ij}$ & Strain tensor & Local deformation from displacement gradients \\
    $\sigma_{ij}$ & Stress tensor & Internal force response \\
    $d,\hat d,\phi$ & Direction quantities & Source--probe directional interpretation \\
    \bottomrule
  \end{tabular}
\end{frame}

\begin{frame}{Definition of the source-probe query and what the prototype estimates}
  \begin{columns}[T]
    \begin{column}{0.58\textwidth}
      \includegraphics[width=\textwidth]{img/web_visualization}
    \end{column}
    \begin{column}{0.38\textwidth}
      \footnotesize
      A source point and a probe point define the directional query.
      \[
        d = x_p - x_s,\qquad
        L = \|d\|
      \]
      \[
        \hat d = \frac{d}{\|d\|},\qquad
        \phi = \operatorname{atan2}(d_y,d_x)
      \]
      The prototype predicts thermal and displacement response together
      with geometric directional quantities such as distance and azimuth.
    \end{column}
  \end{columns}
\end{frame}

\section{Prediction Framework}

\begin{frame}{From Physics to a Prediction Task}
  \begin{block}{Prediction view}
    \[
      \mathcal{F}: X \rightarrow \widehat{Y}
    \]
  \end{block}
  \vskip 4pt
  \begin{columns}[T]
    \begin{column}{0.48\textwidth}
      \begin{block}{Input}
        \footnotesize
        Material parameters, source location, probe location, propagation
        direction, and observation time.
      \end{block}
    \end{column}
    \begin{column}{0.48\textwidth}
      \begin{block}{Output}
        \footnotesize
        Estimated temperature, displacement, directional response,
        and travel time.
      \end{block}
    \end{column}
  \end{columns}
  \vskip 6pt
  The thesis does not claim to solve the complete coupled PDE system
  for every field point. It compares structured predictive estimates.
\end{frame}

\begin{frame}{COMSOL-Based Synthetic Dataset}
  \begin{block}{Controlled thermoelastic experiment}
    2D plane-strain rock domain $1\,\mathrm{m}\times1\,\mathrm{m}$ with a
    heated internal source used to generate time-dependent synthetic data.
  \end{block}
  \vskip 4pt
  \begin{itemize}
    \item Initial rock temperature: $293.15$ K
    \item Source temperature: $1500$ K
    \item Outputs exported per node and time step: $T$, $(u,v)$, velocity,
          stress, strain, and material features
    \item Purpose: create a common reference dataset for model comparison
  \end{itemize}
\end{frame}

\begin{frame}{Data Representation and Preprocessing}
  \begin{itemize}
    \item COMSOL CSV exports are aligned into a single node set.
    \item Dynamic features are assembled into a spatio-temporal tensor:
          \[
            \mathbf{S}\in\mathbb{R}^{N_t \times N_p \times F_d}
          \]
    \item Static material and geometry features are stored separately.
    \item Feature channels are normalized before training:
          \[
            \widetilde{x}=\frac{x-\mu}{\sigma+\varepsilon}
          \]
    \item The same processed dataset supports coordinate, grid, and graph models.
  \end{itemize}
\end{frame}

\section{Comparative Framework}

\begin{frame}{Comparative Prototype Framework}
  \begin{block}{One workflow for all model routes}
    The frontend, backend orchestrator, material catalog, and model
    services use one shared input--output contract for fair comparison.
  \end{block}
  \vskip 6pt
  \begin{itemize}
    \item Same source--probe geometry for all routes
    \item Same effective geological material presets
    \item Same normalized response format: temperature, displacement,
          direction, travel time, and diagnostics
  \end{itemize}
\end{frame}

\begin{frame}{Compared Model Families}
  \footnotesize
  \begin{tabular}{@{}lcccc@{}}
    \toprule
    \textbf{Model} & \textbf{Physics use} & \textbf{Geometry} & \textbf{Strength} & \textbf{Current status} \\
    \midrule
    PINN & Explicit residuals & Coordinates & Physics-informed baseline & Stable \\
    FNO & Supervised operator & Structured grid & Global field patterns & Needs calibration \\
    MeshGraphNet & Supervised rollout & Mesh graph & Mesh-aware prediction & Fixed horizon \\
    Transformer & Supervised operator & Query points & Fast flexible inference & Limited physics priors \\
    \bottomrule
  \end{tabular}
\end{frame}

\begin{frame}{PINN as the Physics-Informed Baseline}
  \begin{block}{Core idea}
    Point-wise MLP
    \[
      f_{\theta}(x,y,t,E,\nu,\rho,\alpha,k,C_p)\rightarrow(T,u,v)
    \]
    trained with supervised loss plus thermoelastic residual terms.
  \end{block}
  \vskip 4pt
  \[
    \mathcal{L}
    =
    \lambda_{\mathrm{sup}}\mathcal{L}_{\mathrm{sup}}
    +
    \lambda_{\mathrm{vel}}\mathcal{L}_{\mathrm{vel}}
    +
    \lambda_{\mathrm{el}}\mathcal{L}_{\mathrm{elastic}}
    +
    \lambda_{\mathrm{th}}\mathcal{L}_{\mathrm{thermal}}
  \]
  \vskip 4pt
  PINN is the only route in the comparison whose objective explicitly
  includes the governing thermoelastic relations.
\end{frame}

\begin{frame}{Other Surrogate Routes}
  \begin{columns}[T]
    \begin{column}{0.32\textwidth}
      \begin{block}{FNO}
        \footnotesize
        Field-to-field neural operator using spectral convolutions on
        structured 2D representations.
      \end{block}
    \end{column}
    \begin{column}{0.32\textwidth}
      \begin{block}{MeshGraphNet}
        \footnotesize
        Mesh nodes become graph nodes; message passing models local
        interactions on irregular geometry.
      \end{block}
    \end{column}
    \begin{column}{0.32\textwidth}
      \begin{block}{Transformer}
        \footnotesize
        Coordinate-query neural operator with tokens carrying geometry,
        state variables, and material features.
      \end{block}
    \end{column}
  \end{columns}
  \vskip 6pt
  \begin{alertblock}{Important framing}
    These models are treated as \emph{surrogate routes for comparison},
    not as physically equivalent solvers.
  \end{alertblock}
\end{frame}

\section{Experimental Results}

\begin{frame}{Calculation Status}
  \begin{block}{Evaluation setup}
    \begin{enumerate}
      \item \textbf{40 scenarios per material}, four materials:
            granite, sandstone, limestone, and basalt.
      \item \textbf{$\mathbf{40\,\times\,4\,\times\,4 = 640}$ model
            responses} evaluated under one shared contract.
      \item \textbf{2D plane-strain only}: the operating envelope of
            all trained surrogates.
    \end{enumerate}
  \end{block}
  \begin{alertblock}{FNO caveat}
    \footnotesize
    FNO is included in the comparison, but its output scales are still
    miscalibrated. It should currently be treated as a diagnostic branch.
  \end{alertblock}
\end{frame}

\begin{frame}{Travel-Time Predictions}
  \begin{columns}[T]
    \begin{column}{0.62\textwidth}
      \includegraphics[width=\textwidth]{img/travel_time_by_model_material}
    \end{column}
    \begin{column}{0.34\textwidth}
      \footnotesize
      \begin{block}{Reading the plot}
        \begin{itemize}
          \item PINN $\approx 0.11$\,ms
          \item Transformer $\approx 0.14$\,ms
          \item FNO $\approx 1.29$\,ms
          \item MeshGraphNet $\approx 2.47$\,ms
        \end{itemize}
      \end{block}
      Model-family differences are larger than the basalt--sandstone
      material difference in the current prototype.
    \end{column}
  \end{columns}
\end{frame}

\begin{frame}{Accuracy Against Analytical Reference}
  \begin{columns}[T]
    \begin{column}{0.60\textwidth}
      \includegraphics[width=\textwidth]{img/travel_time_relative_error}
    \end{column}
    \begin{column}{0.36\textwidth}
      \footnotesize
      \begin{block}{Analytical reference}
      \[
        t_{\mathrm{gt}} = \frac{r}{V_p}
      \]
      \end{block}
      \begin{tabular}{@{}lcc@{}}
        \toprule
        \textbf{Model} & \textbf{Median rel.\ error} & \textbf{Within 10\,\%} \\
        \midrule
        PINN         & $2.0\%$    & $100\%$ \\
        Transformer  & $29.2\%$   & $0\%$   \\
        FNO          & $1403.6\%$ & $0\%$   \\
        MGN          & $> 100\%$  & $0\%$   \\
        \bottomrule
      \end{tabular}
      \vskip 4pt
      PINN is the only route that remains inside the $10\%$ band across
      the balanced sweep.
    \end{column}
  \end{columns}
\end{frame}

\begin{frame}{Inference Latency and Time Behavior}
  \begin{columns}[T]
    \begin{column}{0.62\textwidth}
      \includegraphics[width=\textwidth]{img/inference_latency_by_model}
    \end{column}
    \begin{column}{0.34\textwidth}
      \footnotesize
      \begin{block}{Median latency}
        \begin{itemize}
          \item Transformer $\approx 1.3$\,ms
          \item MeshGraphNet $\approx 1.5$\,ms
          \item FNO $\approx 2.4$\,ms
          \item PINN $\approx 2.6$\,ms
        \end{itemize}
      \end{block}
      \vskip 4pt
      PINN is also the only route that changes smoothly with requested
      observation time; the other routes currently behave as fixed-horizon
      predictors.
    \end{column}
  \end{columns}
\end{frame}

\begin{frame}{Threats to Validity}
  \begin{enumerate}
    \item \textbf{Limited training data.}
          The COMSOL-derived pilot configuration is not a full calibrated
          geological database.
    \item \textbf{Simplified 2D assumptions.}
          The study uses effective material parameters in a plane-strain setup.
    \item \textbf{FNO normalization issue.}
          The current 2D adaptation is not yet physically reliable.
    \item \textbf{Analytical reference is isotropic.}
          $r/V_p$ is a useful sanity check, not a complete thermoelastic validation.
  \end{enumerate}
\end{frame}

\section{Conclusion}

\begin{frame}{Conclusions}
  \begin{enumerate}
    \item The thesis develops a physically grounded
          \emph{comparative framework} for AI-guided directional prediction.
    \item PINN is the strongest \emph{physics-informed baseline} in the
          current comparison and shows the best agreement with the analytical
          travel-time reference.
    \item Transformer is the fastest route, while MeshGraphNet remains a
          useful mesh-aware surrogate candidate.
    \item FNO should currently be treated as a diagnostic branch until
          output normalization is corrected.
  \end{enumerate}
\end{frame}

\begin{frame}{Future Work}
  \begin{enumerate}
    \item Retrain and recalibrate the 2D FNO route.
    \item Add time-conditioned rollout behavior to MeshGraphNet.
    \item Extend the material catalog with more complete thermoelastic
          parameters and validation data.
    \item Compare against additional analytical or numerical baselines.
  \end{enumerate}
  \vskip 8pt
  \begin{exampleblock}{Final message}
    The prototype estimates directional characteristics of thermoelastic
    wave propagation within the simplified assumptions of the study.
  \end{exampleblock}
\end{frame}

\begin{frame}{Thank You}
  \centering
  \Large Questions?
\end{frame}

\end{document}
